%% kapa: exam=1 round=26 year=2015 period=1 form=A subject=economics subject_ko=경제원론 seq=16 number=56
%% source: Q-Net 감정평가사 기출문제 (cst00309) artlSeq=5138001 file=extracted/1008/5138001/제26회_감정평가사 1차_기출/제26회_감정평가사 1차_기출/1교시_감평_A형.pdf
%% conversion: kordoc 4.24.51 -> tools/gen1.py (auto); review=proofread-sonnet; flags=none
\begin{question}{56}
생산측면에서 외부효과가 발생하는 경우에 관한 설명으로 옳지 않은 것은?
\begin{choices}
  \item 부정적 외부효과가 존재할 경우, 시장균형거래량에서 사회적 한계비용이 시장균형가격보다 낮다.
  \item 긍정적 외부효과가 존재할 경우, 시장균형거래량은 사회적 최적거래량보다 작다.
  \item 부정적 외부효과가 존재할 경우, 경제적 순손실(자중손실)이 발생한다.
  \item 긍정적 외부효과가 존재할 경우, 경제적 순손실(자중손실)이 발생한다.
  \item 외부효과는 한 사람의 행위가 제3자의 경제적 후생에 영향을 미치고, 그에 대한 보상이 이루어지지 않을 때 발생한다.
\end{choices}
\end{question}
\answer{1}
